\documentclass[11pt]{article}
\usepackage{icmlko}
\begin{document}

\title{다섯 행이 여덟 노트를 만들었다 --- ``축 개수''는 어느 수집이 붙었냐의 다른 이름}
\author{월드모델 실험실}
\date{노트 488 · 2026년 8월 3일}
\maketitle

\begin{abstract}
노트 487 이 모바일에서 저덮음 축을 \emph{한꺼번에} 가리면 축 많은 절반이
$+0.0167$ 오른다는 것을 보고 ``어느 축인가''를 남겼다. 갈랐더니 이름이
나왔다 --- 그리고 그 이름이 노트 $480 \sim 487$ 여덟 노트의 계기를 통째로
무너뜨린다.

\begin{itemize}
\item \textbf{많은 쪽의 $+0.0167$ 은 \texttt{media\_push} 하나다}
($+0.0193$, 몫 $116\%$). 그 축은 모바일 유보 $441$ 행 중
\textbf{다섯 행}에만 있고, 그 다섯이 전부 많은 쪽에 있다.
\item \textbf{적은 쪽의 $-0.0148$ 은 다른 축이다} --- \texttt{gen}
($-0.0150$, 몫 $101\%$). \textbf{두 절반이 서로 다른 축에 매여 있다.}
\item \textbf{많은 쪽은 $28$ 행이다.} 모바일 유보 $441$ 행 중 $375$ 행이
축 개수 $15$ 로 \emph{똑같다}. 중앙값 나누기가 성립하지 않는다.
\item \textbf{모바일만이 아니다} --- 다섯 도메인 중 \textbf{넷}이 제일
흔한 축 개수 하나에 절반 넘게 뭉쳐 있다(만화 $71\%$ · 애니 $72\%$ ·
세계애니 $62\%$ · 모바일 $85\%$). 퍼진 것은 게임($35\%$)뿐이다.
\end{itemize}

\textbf{그래서 ``축이 많은 행''은 대개 ``축이 한 개 더 있는 행''이고, 그
한 개는 늘 특정 수집이다.}
\end{abstract}

\begin{figure}[t]
\centering
\includegraphics[width=\textwidth]{figs/f1.pdf}
\caption{\textbf{가} --- 모바일 축 개수 분포. \textbf{나} --- 살아 있는
저덮음 축별 분해. \textbf{다} --- 다섯 도메인의 뭉침.}
\end{figure}

\section{계기가 무엇이었나}

노트 $480 \sim 487$ 은 ``축이 많은 행과 적은 행의 격차''를 독립변수로
썼다. 실제로 그 변수가 무엇이었는지 각 도메인에서 두 절반의 축 덮음을
견주면 나온다.

\begin{center}
\begin{tabular}{llrr}
\hline
도메인 & 반을 가른 축 & 많은 쪽 & 적은 쪽 \\
\hline
애니 & \texttt{venue\_prominence} & $100\%$ & $2\%$ \\
모바일 & \texttt{wiki\_*} 셋 & $82\%$ & $0\%$ \\
만화 & \texttt{goods\_scale} & $79\%$ & $1\%$ \\
세계애니 & \texttt{wiki\_*} 셋 & $100\%$ & $36\%$ \\
게임 & \texttt{media\_push} & $85\%$ & $7\%$ \\
\hline
\end{tabular}
\end{center}

\textbf{도메인마다 수집 하나의 있음/없음이 반을 갈랐다.} 애니는 사실상
\texttt{venue\_prominence} 표시자고($100\%$ 대 $2\%$), 모바일은 위키가
붙었냐다.

이것이 노트 486 을 뒤에서 설명한다. 거기서 ``\texttt{venue\_prominence}
있는 행은 잘 기록된 부분 무리이고 라벨이 눌려 있다''고 했는데, 그
부분 무리가 바로 \emph{애니의 축 많은 절반 그 자체}였다. 두 발견은 같은
것을 두 번 말한 것이다.

\section{저덮음의 대부분은 `저'가 아니라 `없음'이다}

부수적으로 하나 더. 도메인별 ``저덮음 축''을 세면 $22 \sim 32$ 개인데,
\textbf{그 대부분이 덮음 $0\%$} --- 이름표만 있고 그 도메인 유보에는 한
행도 없다. 실제로 살아 있는 저덮음 축은 게임 $7$ · 모바일 $5$ · 만화 $4$
· 세계애니 $4$ · \textbf{애니 $2$} 개뿐이다.

\begin{quote}
\textbf{조항} --- \textbf{덮음 $0\%$ 축과 덮음 $10\%$ 축을 한 통에 넣지
않는다.} 앞엣것은 가려도 아무 일이 없고(정확히 $+0.0000$), 뒤엣것은
몇 행을 움직인다. ``저덮음 축 $N$ 개''라고 쓸 때 $N$ 은 거의 늘 부풀려진
수다.
\end{quote}

\section{스스로 걸린 것}

첫 실행은 낱개 가림을 덮음 $95\%$ \emph{이상} 축에 돌리고 ``전체''는
$95\%$ \emph{미만} 축을 가렸다. 두 집합이 겹치지 않으니 ``합
$-0.2482$ 대 전체 $+0.0167$''이 나왔고, 하마터면 \emph{나무의
비선형성}으로 읽을 뻔했다. 표의 축이 전부 덮음 $100\%$ 인 것을 보고
잡았다.

\begin{quote}
\textbf{조항} --- \textbf{분해는 합이 전체와 맞는지 먼저 본다.} 안 맞으면
비선형을 말하기 전에 \emph{같은 것을 쪼갠 게 맞는지} 확인한다. 고친 뒤
모바일은 합 $+0.0257$ 대 전체 $+0.0167$ 로 맞았다.
\end{quote}

\section{여덟 노트를 다시 쓴다}

노트 487 이 ``한 물음에 답이 넷''이라고 닫았다. 이 노트는 그보다
간단하게 닫는다.

\begin{quote}
\textbf{격차는 축의 \emph{개수}가 아니라 어느 \emph{수집}이 붙었냐다.}
\end{quote}

그리고 모바일의 ``이상''은 --- 여덟 노트를 끈 그것은 --- \textbf{$441$ 행
중 다섯 행에 붙은 축 하나}였다. $28$ 행짜리 절반 위에 서 있었다.

\textbf{모형 이야기가 아니었다. 처음부터 수집 이야기였다.}

\section*{남은 것}

\begin{itemize}
\item \textbf{게임은 진짜로 퍼져 있다}($35\%$, $85$ 대 $95$ 행). 유일하게
성립하는 나누기이므로 노트 $480 \sim 487$ 의 물음을 \textbf{게임에서만}
다시 물으면 답이 남는지 볼 수 있다.
\item \textbf{그 다섯 행이 무엇인가} --- \texttt{media\_push} 가 붙은
모바일 다섯 행을 직접 보면 왜 해로운지 알 수 있다.
\item \textbf{조항 하나 더 필요하다} --- 이런 반 가르기를 쓰기 전에
\emph{분포부터 찍는} 가드. 넷이 뭉쳐 있었는데 여덟 노트 동안 아무도 안
봤다.
\end{itemize}

\end{document}
